\chapter{Listes de lecture}\label{chap:listes}

\section{Plusieurs listes}
Vous pouvez créer autant de listes que vous voulez. Deux existent d'emblée : la liste par défaut, qui reprend votre
ancienne liste lors de la mise à jour, et les favoris, que vous remplissez avec le bouton $\equiv$+ de l'écran
Lecture.

La liste affichée est aussi celle qui se lit. Quand vous changez de liste avec le sélecteur, nous enregistrons la
liste que vous quittez et chargeons la nouvelle ; le morceau en cours continue, et le bouton « suivant » suivra la
nouvelle liste. Nous avons fait ce choix pour que l'écran montre toujours ce que vous allez entendre.

\capture{playlist-en-top}{Le sélecteur de liste}

\section{Gérer vos listes}
Le menu du sélecteur permet de créer une liste, de renommer ou dupliquer la liste affichée, de la supprimer (une
autre liste prend alors sa place) et d'échanger des listes au format M3U.

\section{Ajouter des morceaux}
Depuis l'écran Lecture, le bouton $\equiv$+ puis « Add to playlist\dots » range le morceau en cours ; vous pouvez
aussi créer une liste au passage. Depuis le navigateur, sélectionnez des fichiers ou des dossiers puis touchez
« Add ». Depuis une liste, sélectionnez des morceaux et copiez-les dans une autre liste.

\section{La file d'attente}\label{sec:file}
« Play next » insère des morceaux juste après celui en cours, sans toucher à vos listes. Ils passent dans l'ordre où
vous les avez ajoutés, puis la lecture reprend là où elle en était. Si vous choisissez vous-même un autre morceau
entre-temps, la file est conservée mais ce point de reprise est oublié.

\section{Import et export M3U}
L'export écrit la liste affichée au format M3U étendu, avec titre, auteur et durée. Les emplacements sont complets,
y compris à l'intérieur des archives et des catalogues en ligne : une liste exportée se réimporte à l'identique.
L'import accepte aussi les listes d'autres lecteurs et crée une nouvelle liste portant le nom du fichier. Vos
anciennes listes au format XSPF restent accessibles dans le navigateur.

\section{Trier et ranger}
Le tri porte sur le titre, l'auteur ou la durée. Pour un ordre à votre goût, déplacez les morceaux par leur poignée.
Les statistiques donnent le nombre de morceaux, le nombre de fichiers distincts et la durée totale.
